\begin{abstract}
Federated learning (FL) keeps training data on client devices, but the network
traffic it generates can still reveal which model is being trained. We present
CHAI-FP, a framework that characterizes how much AI model identity a passive
observer can infer from packet-level metadata alone, without inspecting payloads
or model parameters. CHAI-FP reconstructs the round structure of FL traffic and
identifies models hierarchically at the family, architecture, and variant
levels. Its key element is cross-device fingerprinting. Because FL clients train the
same model on heterogeneous hardware and operating systems, the differences
between their traffic reveal the architecture's compute profile. On a physical testbed with three
heterogeneous clients (two Ubuntu, one Windows) and 32 experiments spanning convolutional, recurrent,
transformer, and autoencoder models, cross-device fingerprinting identifies the
family of unseen models with a balanced accuracy of 0.74, compared with 0.58 for
per-client classification, and raises family accuracy in 5~s windows from 0.31
to 0.51. Within a family, the architecture is identified with a mean balanced
accuracy of 0.79 from complete traces, and convolutional architectures with 0.71
from 5~s windows. The model payload size is a baseline leak that reveals the
exact variant in every trace, and in 69\% from a single model download within
a 5~s window, but it is removed by padding, whereas family identification from
timing alone remains at 0.62, showing that concealing model
identity in FL requires timing defenses in addition to size padding.
\end{abstract}

\begin{IEEEkeywords}
Federated learning, traffic analysis, side channels, model fingerprinting,
network security, privacy.
\end{IEEEkeywords}

\section{Introduction}
\label{sec:introduction}
Artificial intelligence (AI) workloads are increasingly executed across edge
devices, cloud services, and distributed learning
systems~\cite{zhou2019edge,mcmahan2017communication}. Although encryption and
access controls protect transmitted content and model artifacts, they do not
conceal all characteristics of the communication generated during execution:
packet sizes, directions, timing, volumes, and burst patterns remain observable
to a network intermediary~\cite{sirinam2018deep}. Such metadata may reveal
properties of an AI workload even when the observer cannot inspect model
parameters, application data, or payloads~\cite{shuvo2025fingerprinting}.
Federated learning (FL) is particularly exposed, because every training round
produces a characteristic exchange in which each client downloads the global
model, trains locally, and uploads an update.

\textbf{Existing Research:} Prior studies infer neural network information from
indirect execution observations, including memory accesses, shared caches, GPU
execution traces, and system resource
measurements~\cite{hua2018reverse,yan2020cache,hu2020deepsniffer,patwari2022dnn}.
These observations require physical access to, or co-location with, the
executing device. Network traffic is a more remote observation surface, on which
communication protocols, training schedules, transport behavior, and network
conditions may obscure architectural differences. Encrypted-traffic analysis
reveals visited
websites~\cite{panchenko2011website,wang2014effective,hayes2016kfingerprinting,sirinam2018deep},
and PACKETPRINT~\cite{li2022packet} identifies mobile applications and user
activities from encrypted wireless packet sequences, although its hierarchy
captures traffic patterns rather than structural relationships among AI models.
For AI workloads, network traffic has been used to discriminate CNN and RNN
families in FL~\cite{shuvo2025fingerprinting}, to classify CNN and RNN
architectures and variants from encrypted wireless traffic with results reported
primarily at the family level~\cite{shuvo2026flare}, to identify AI applications
and their execution phases~\cite{oyaniyi2026fingerprinting}, and to identify AI
families from 5G scheduling metadata~\cite{shuvo2026flint}.

\textbf{Gap:} The extent to which passive network observations reveal the
structural identity of AI models beyond the family level remains insufficiently
characterized. It is not established whether an observer can progressively
distinguish families, architectures, and variants, how much of this information
is available within short observation windows, or whether it survives
countermeasures such as size padding. Moreover, existing approaches classify each
client's traffic independently and do not exploit a defining property of FL:
all clients train the same model concurrently on heterogeneous hardware and
operating systems.

\textbf{Motivation:} Three technical challenges motivate this investigation.
First, different families and architectures may generate similar communication
patterns under common FL protocols, so network similarity is an imperfect
indicator of model structure. Second, discriminative characteristics of complete
traces may be unavailable within short windows, which capture only portions of a
round. Third, observations from the same execution are correlated, so randomly
distributing them between training and testing produces optimistic estimates;
generalization must instead be assessed on models absent from training.
Architectural leakage also has security consequences. Knowledge of a model's
family, architecture, or variant narrows the range of potential targets and
helps an adversary select relevant attack techniques, for instance by training a
surrogate of the matching architecture for model extraction or transfer-based
adversarial attacks~\cite{TODO-model-extraction}, and it exposes proprietary
design choices.

\textbf{This Work:} We present CHAI-FP, a conditional hierarchical framework
that infers the family, architecture, and variant of the model under training
from passive packet metadata. CHAI-FP reconstructs the round structure of FL
traffic, from which it derives communication features and a passive estimate of
local training time, the compute gap between model download and upload, and it
combines the traffic of all clients into a cross-device representation. We
evaluate it on a physical FL testbed with three clients that differ in both
hardware and operating system, and 32
experiments, using model-disjoint validation on 0.5 to 5~s windows and complete
traces.

\noindent\textbf{Contributions:} The contributions of this study are threefold.
(i)~A conditional hierarchical fingerprinting framework that identifies model
families, architectures, and variants from passive network metadata, combining
packet-morphology features with compute-timing features derived from the FL
round structure; within a family, packet-morphology features identify the
architecture with a mean balanced accuracy of 0.79.
(ii)~Cross-device fingerprinting, which exploits the heterogeneous clients of a
federated job to raise family identification of unseen models from 0.58 to 0.74
for complete traces and from 0.31 to 0.51 for 5~s windows, where it also
identifies convolutional architectures with 0.71; an ablation shows that the gain requires both
the cross-device representation and a regularized linear classifier.
(iii)~A model-disjoint evaluation of the observable leakage, showing that
payload size is a baseline leak that reveals the exact variant in every trace,
and in $69\%$ from a single model download within a 5~s window, but is
removed by padding, whereas family identification from timing
alone remains at 0.62, so that concealing model identity requires timing
defenses.

The remainder of the paper is organized as follows.
Section~\ref{sec:setup} describes the experimental setup,
Section~\ref{sec:methodology} presents the CHAI-FP framework,
Section~\ref{sec:results} reports the results,
Section~\ref{sec:discussion} discusses countermeasures, and
Section~\ref{sec:conclusion} concludes the paper.
